% !TEX root = ../main_ext_1.tex

\section*{Abstract}
\thispagestyle{empty}

Vector autoregressions (VARs) and local projections (LPs) estimate the same population impulse
responses, yet they are usually compared at a common lag order, which confounds the choice of
estimator with the choice of model complexity. We revisit this comparison in a Monte Carlo study
that holds effective complexity constant, following the complexity-adjusted benchmark of
\textcite{ludwig2026local}. Under correct specification, a fixed LP(4) already behaves like a
highly parameterized VAR of order $q \approx 15$, so the familiar point-estimation edge of
low-order VARs reflects parsimony rather than an intrinsic advantage of the iterated method. We
then introduce dynamic misspecification through a VARMA(4,1) process. This does not overturn the
point-estimation ranking, as the equal-lag VAR(4) keeps its root mean squared error advantage over
LP(4) at every level of persistence and misspecification. The decisive difference is inferential.
Standard VAR confidence intervals understate the sampling dispersion and undercover badly near the
unit-root boundary, while the heteroskedasticity-robust intervals of the local projection stay near
nominal. Population equivalence of the two estimands therefore does not carry over to finite-sample
inference, and the better method depends on whether the researcher prioritizes the efficiency of the
VAR point estimate or the interval validity of the direct projection.

\clearpage
